%
% 12-partial.tex -- Stammfunktionen für Partialbrüche
%
% (c) 2026 Prof Dr Andreas Müller
%
\subsection{Stammfunktionen für Partialbrüche}
Dank der Partialbruchzerlegung zerfällt in Summanden der Form
\[
\frac{A}{(z-\alpha)^l}
=
A(z-\alpha)^{-l}.
\]
Die Stammfunktion ist
\[
\int \frac{A}{(z-\alpha)^l} \,dz
=
\begin{cases}
\displaystyle
-\frac{A}{(l-1)(z-\alpha)^{l-1}}
&\qquad\text{für $l>1$}
\\[10pt]
A\log(z-\alpha)
&\qquad\text{für $l=1$.}
\end{cases}
\]
Diese Stammfunktionen erlauben eine Stammfunktion für die
Funktion $f(z)$ zu finden.
Sie besteht aus zwei Teilen.
Ein Teil hat die Form einer rationalen Funktion, die aus den
Beiträgen des Polynoms $p(z)$ und den Beiträgen der
Bruchterme $A_k^l/(z-\alpha_k)^l$ mit $l>1$ gelten.
Dazu kommen als zweiter Teil die Funktionen
\[
\sum_{k=1}^n A_k^1\log(z-\alpha_k).
\]
Um eine Stammfunktion eines Integranden im Körper der rationalen Funktionen
zu konstruieren, muss man den Funktionenkörper um Logarithmen der
Linearfaktoren des Nenners erweitern.

